This appendix reports additional analyses that assess the robustness and
interpretation of the main capability--drift association, rather than to
establish causal effects.
All analyses use the frozen Normal-framework records described in the project
repository.  Missing rollouts are retained when available, and paired analyses
use only the common evaluable samples for the relevant model pair.

\subsection{Paired Within-Sample Comparisons}
\label{app:additional-paired}

To avoid treating model-level averages as independent observations, we compare
models within the same samples.  For each model and sample, we compute the mean
drift indicator over the available rollouts and test the paired differences with
the two-sided Wilcoxon signed-rank test.  The confidence intervals are percentile
bootstrap intervals over samples (10{,}000 resamples); the four reported
comparisons use Holm correction for the Wilcoxon $p$-values.

\begin{table}[t]
\centering
\small
\setlength{\tabcolsep}{3pt}
\resizebox{\columnwidth}{!}{%
\begin{tabular}{l r r r r r}
\toprule
Comparison & $\Delta$DR & $r_{\rm rb}$ & $p_{\rm Holm}$ & 95\% CI & $N$ \\
\midrule
Qwen3.6-35B $-$ Qwen3.5-4B & $+0.154$ & $+0.623$ & $2.44\times10^{-20}$ & $[+.124,+.184]$ & 629 \\
Mistral-24B $-$ Ministral-8B & $-0.074$ & $-0.307$ & $2.72\times10^{-6}$ & $[-.103,-.043]$ & 641 \\
Qwen3.6-35B $-$ Claude Sonnet 4.6 & $+0.036$ & $+0.169$ & $0.0167$ & $[+.007,+.066]$ & 632 \\
Llama-4 Maverick $-$ Llama-3.1-8B & $-0.094$ & $-0.730$ & $1.65\times10^{-15}$ & $[-.116,-.073]$ & 638 \\
\bottomrule
\end{tabular}
}%
\caption{Paired comparisons of per-sample mean drift. $\Delta$DR is model A
minus model B; $r_{\rm rb}$ is the rank-biserial effect size. The analysis
supports heterogeneous within-family directions, but does not identify the
causal source of those differences.}
\label{tab:additional-paired}
\end{table}

The confidence intervals exclude zero for all four pre-specified comparisons,
including both the increasing Qwen contrast and the decreasing Mistral contrast.
This supports the narrower claim that parameter count alone is insufficient to
predict authorization drift; it does not by itself establish that alignment
strategy is the causal explanation.

\subsection{Correlation and Covariate Sensitivity}
\label{app:additional-correlation}

Across the 13 paper models, the model-level TSR--DR association is positive
(Pearson $r=0.858$, $p<0.001$; Spearman $\rho=0.929$, $p<0.001$).  Leave-one-
model-out Pearson correlations range from $0.799$ to $0.894$.  These are
descriptive model-level statistics and should not be interpreted as independent
replicates or as causal evidence.

As a sensitivity analysis, we fit trajectory-level binomial GLMs with
sample-clustered standard errors.  The coefficient of task success remains
positive after adding family and domain indicators (OR $=1.023$, $p=0.823$),
and is larger but imprecisely estimated when tool-call count is added (OR
$=1.221$, $p=0.057$).  Adding annotated trigger count gives OR $=1.025$
($p=0.805$), while the fuller model with model, domain, tool-count, and
trigger-count terms gives OR $=1.156$ ($p=0.178$).  The wide uncertainty in
these adjusted specifications, and the possibility that tool-call count is a
mediator rather than a confounder, make them sensitivity analyses rather than
adjusted causal estimates.

\begin{table}[t]
\centering
\small
\setlength{\tabcolsep}{3pt}
\resizebox{\columnwidth}{!}{%
\begin{tabular}{l r r l}
\toprule
Specification & OR & $p$ & 95\% CI \\
\midrule
Family + domain & 1.023 & .823 & $[.841,1.243]$ \\
Family + domain + tools & 1.221 & .057 & $[.994,1.499]$ \\
Family + domain + triggers & 1.025 & .805 & $[.841,1.251]$ \\
Model + domain + tools + triggers & 1.156 & .178 & $[.936,1.428]$ \\
\bottomrule
\end{tabular}
}%
\caption{Trajectory-level sensitivity models for the success--drift association.
All models use sample-clustered standard errors.}
\label{tab:additional-glm}
\end{table}

\subsection{Exploratory Behavioral Annotation of First Drift}
\label{app:additional-behavior}

To probe whether first drift actions are behaviorally related to task progress,
we sampled 600 drifting trajectories, 200 each from Claude Sonnet 4.6,
Qwen3.5-397B-A17B, and Qwen3.6-27B.  Three independent LLM judges assessed
task advancement and competing error evidence from the observable trajectory,
without being shown the source-agent identity or the pre-existing drift label.
Consensus was formed by a two-of-three vote; all-different cases were treated
as ambiguous.  The resulting annotation is exploratory and is not human gold
annotation, evidence of hidden intent, or a causal test.

At the behavioral level, 85.4\% of first violations were classified as
instrumental without competing error evidence and a further 12.0\% as
instrumental with error evidence.  Thus 97.4\% were judged task-instrumental
under this operational rubric.  However, raw task-relation agreement was
92.4\%, while Fleiss' $\kappa$ was only 0.045 (Gwet's AC1 $=0.921$), reflecting
the highly imbalanced label distribution.  The complementary failure-reason
annotation had Fleiss' $\kappa=0.085$ and should therefore be read as a
behavioral sensitivity analysis rather than a validated decomposition of model
intent.

\begin{table}[t]
\centering
\small
\setlength{\tabcolsep}{3pt}
\resizebox{\columnwidth}{!}{%
\begin{tabular}{l r r r r}
\toprule
Behavioral pattern & Claude & Qwen-397B & Qwen-27B & Equal-weight \\
\midrule
Instrumental, no competing error & 74.9 & 91.4 & 89.8 & 85.4 \\
Instrumental, error evidence & 22.0 & 6.6 & 7.3 & 12.0 \\
Non-instrumental failure & 0.8 & 0.4 & 0.0 & 0.4 \\
Anomalous behavior & 0.6 & 1.1 & 0.0 & 0.6 \\
Ambiguous / unresolved & 1.6 & 0.4 & 2.9 & 1.7 \\
\bottomrule
\end{tabular}
}%
\caption{Exploratory three-judge annotation of first violating calls. Entries
are percentages; the source-model columns use sample-clustered bootstrap 95\%
confidence intervals in the machine-readable supplement.}
\label{tab:additional-behavior}
\end{table}

The study does not replace human validation, and no human--judge agreement is
reported for it; the human quality audit of the benchmark samples themselves is
reported separately in Appendix~\ref{app:human-eval}.
